% \section{Related Work}
% \label{sec:related_work}

% Due to space constraints, we defer most Related Work to Appendix~\ref{app:sec:related_work} and focus here on prior research most relevant to our contribution.
% The precise details of \citet{hoffman2022chinchilla} have recently come under scrutiny, leading to a number of important replication and re-evaluation studies.
% For instance, Chinchilla used three different approaches, two of which agreed with each other, but the third did not; \citet{besiroglu2024chinchillascalingreplicationattempt} conducted a detailed investigation of this third analysis and found that it could be made consistent with the first two analyses by fixing optimizer issues and not rounding reported fit parameters.
% In a similar vein, \citet{porian2024resolvingdiscrepancies} and \citet{pearce2024reconciling} sought to resolve a discrepancy between \citet{hoffman2022chinchilla} and \citet{kaplan2020scaling} on how to scale data and parameters to produce the best performing model; \citet{porian2024resolvingdiscrepancies} found that the discrepancy could be resolved by three differences (last layer computational cost, warmup duration, and scale-dependent optimizer tuning) while \citet{pearce2024reconciling} found that much of the discrepancy could be attributed to \citet{kaplan2020scaling} counting only non-embedding parameters.

% Like \citet{besiroglu2024chinchillascalingreplicationattempt} and \citet{porian2024resolvingdiscrepancies}, our work scrutinizes the seminal work of Chinchilla.
% However, our analyses focuses specifically on how robust the original Chinchilla methodology and results are to different perturbations.
% Our contribution concludes with a direct confirmation of the original findings, providing evidence that Chinchilla's compute-optimal guidance is robust.

% \section{Discussion}
% \label{sec:discussion}

% This work began with a perhaps surprising result: three different interpretations of Chinchilla's model parameters are possible, with discrepancies as high as 15.2\%, but all three support (or strengthen) key Chinchilla results.
% Neither the estimated scaling law parameters  nor the widely adopted ``20-to-1'' compute-optimal tokens-to-parameter ratio changed meaningfully. 
% Indeed, our refitting using the standard formula model parameters suggests an even more stable relationship, with the token-to-parameter ratio varying even less across different pretraining compute budgets.

% To understand this robustness more deeply, we systematically investigated how various hypothetical perturbations would affect key Chinchilla results.
% We perturbed the parameter counts in four structured ways and re-ran the fitting analysis for each. This stress test revealed the specific ways in which different types of errors impact the scaling law parameters. A simple multiplicative error, for example, exponentially shifts the constant in the optimal tokens-per-parameter ratio, while an additive error or a systematic bias can more dramatically alter its trend with respect to the target training compute budget.

% Ultimately, our findings serve as both a critical re-examination and a powerful confirmation of the original Chinchilla results.
% Our subsequent analyses should give practitioners even greater confidence in Chinchilla's compute-optimal prescription.
% Its guidance withstands not only the specific interpretation used, but also a range of other potential perturbations, reinforcing its value as a durable and practical blueprint for the field.

% \paragraph{Future Directions} One obvious next step is to evaluate the robustness of more recent scaling results with additional considerations such as inference constraints \citep{sardana2024beyond}, data constraints \citep{muennighoff2023scaling} and overtraining \citep{gadre2024scalinglawovertraining}.

% % \subsubsection*{Author Contributions}
% % If you'd like to, you may include  a section for author contributions as is done
% % in many journals. This is optional and at the discretion of the authors.

% % \subsubsection*{Acknowledgments}
% % Use unnumbered third level headings for the acknowledgments. All
% % acknowledgments, including those to funding agencies, go at the end of the paper.

\section{Related Work}
\label{sec:related_work}

% \paragraph{Scaling Laws in Neural (Language) Models} 
While initial research on scaling laws in neural models began decades ago \citep{barkai1993scaling, mhaskar1996neural, pinkus1999approximation}, advances in scaling large language models brought such interest into renewed focus \citep{hestness2017deeplearningscalingpredictable, kaplan2020scaling, brown2020language}, causing an explosion of research.
For a non-exhaustive list, theoretical understanding of scaling laws has advanced substantially \citep{spigler2020asymptoticlearningcurves, bousquet2020theoryuniversallearning, hutter2021learningcurvetheory, sharma2022scaling, maloney2022solvable, roberts2022principles, bahri2024explaining, paquette2024fourplus3phases, atanasov2024scaling, bordelon2024dynamical, bordelon2024feature, lin2024scaling, brill2024neural}, complemented by empirical studies \citep{rosenfeld2020constructive, henighan2020scaling, gordon2021dataparamscalingnmt, tay2021scaleefficiently, ghorbani2021scaling, zhai2022scaling, alabdulmohsin2022revisitingscalinglaws, dehghani2023scaling, bachmann2023scalingmlpstaleinductive,everett2024scaling, qiu2025scalingcollapse}.

The precise details of \citet{hoffman2022chinchilla} have recently come under scrutiny, leading to a number of important replication and re-evaluation studies.
For instance, Chinchilla used three different approaches, two of which agreed with each other, but the third did not; \citet{besiroglu2024chinchillascalingreplicationattempt} conducted a detailed investigation of this third analysis and found that it could be made consistent with the first two analyses by fixing optimizer issues and not rounding reported fit parameters.
In a similar vein, \citet{porian2024resolvingdiscrepancies} and \citet{pearce2024reconciling} sought to resolve a discrepancy between \citet{hoffman2022chinchilla} and \citet{kaplan2020scaling} on how to scale data and parameters to produce the best performing model; \citet{porian2024resolvingdiscrepancies} found that the discrepancy could be resolved by three differences (last layer computational cost, warmup duration, and scale-dependent optimizer tuning) while \citet{pearce2024reconciling} found that much of the discrepancy could be attributed to \citet{kaplan2020scaling} counting only non-embedding parameters.

Additional research has also studied how scaling interacts with specific considerations such as efficient inference \citep{sardana2024beyond,bian2025scalinginferenceefficient}, transfer \citep{hernandez2021scalinglawstransfer,barnett2024empiricalstudyscalinglaws}, data quality and diversity \citep{chen2025revisiting,
hernandez2022scalingrepeatdata, muennighoff2023scaling, qin2025scalingsynthetic, shukor2025scalinglawsnativemultimodal}, overtraining \citep{gadre2024scalinglawovertraining}, quantization and precision \citep{dettmers20234bitprecisionscaling,sun2025scalinglawsfloatingpoint,kumar2024scalinglawprecision}, differential privacy \citep{mckenna2025scaling}, distillation \citep{busbridge2025distillationscalinglaws}, model architecture \citep{clark2022unifiedscalinglawsroutedlanguagemodels, kudugunta2023matformer, abnar2025parameters, ludziejewski2025joint, liew2025scaling}, context length \citep{xiong2023effectivelongcontextscalingfoundation,agarwal2024manyshot, arora2024bayesianscalinglawsincontext}, vocabulary size \citep{tao2024scalingvocabulary}, robustness to jailbreaking \citep{howe2025scalingrobustness,anil2024manyshot, hughes2024bestofnjailbreaking}, pruning \citep{rosenfeld2021pruningacrossscales}, multimodality \citep{aghajanyan2023scalinggenerativemultimodallm, cherti2023scalingcontrastivelanguageimagelearning}, fine-tuning \citep{kalajdzievski2024scalinglawsforgettingfinetuning, zhang2024scalingmeetsllmfinetuning} and agents and world models \citep{pearce2025scaling}.

Recent work has also highlighted novel scaling phenomena such as inverse scaling \citep{mckenzie2023inversescalingbiggerisnt,gema2025inversescalingtesttimecompute}, emergent capabilities \citep{srivastava2023imitationgamequantifyingextrapolating, wei2022emergentabilitieslargelanguage, schaeffer2023emergent, hu2024predictingemergentabilitiesinfinite, schaeffer2024elusive, snell2024predictingemergentcapabilitiesfinetuning, wu2024ushapedinverteduscalingemergent}, and critical issues like data contamination \citep{schaeffer2023pretrainingtestsetneed, jiang2024investigatingdatacontaminationpretraining, dominguezolmedo2024trainingtesttaskconfounds} and model-data feedback loops \citep{dohmatob2024taletailsmodelcollapse, gerstgrasser2024modelcollapseinevitablebreaking, kazdan2024collapsethriveperilspromises}.
The advent of so-called ``thinking'' or ``reasoning'' models \citep{jaech2024openai} has sparked a new wave of interest in scaling inference compute \citep{brown2024largelanguagemonkeysscaling, snell2024scaling, wu2024inferencescalinglawsempirical, chen2024simpleprovablescalinglaw, sadhukhan2025kineticsrethinkingtesttimescaling,levi2025simple,schaeffer2025monkeypowerlaws,kwok2025robomonkey}.

Like \citet{besiroglu2024chinchillascalingreplicationattempt} and \citet{porian2024resolvingdiscrepancies}, our work scrutinizes the seminal work of Chinchilla.
However, our analyses focuses specifically on how robust the original Chinchilla methodology and results are to different perturbations.
Our contribution concludes with a direct confirmation of the original findings, providing evidence that Chinchilla's compute-optimal guidance is robust.

\section{Discussion}
\label{sec:discussion}

This work began with a perhaps surprising result: three different interpretations of Chinchilla's model parameters are possible, with discrepancies as high as 15.2\%, but all three support (or strengthen) key Chinchilla results.
Neither the estimated scaling law parameters  nor the widely adopted ``20-to-1'' compute-optimal tokens-to-parameter ratio changed meaningfully. 
Indeed, our refitting using the standard formula model parameters suggests an even more stable relationship, with the token-to-parameter ratio varying even less across different pretraining compute budgets.

To understand this robustness more deeply, we systematically investigated how various hypothetical perturbations would affect key Chinchilla results.
We perturbed the parameter counts in four structured ways and re-ran the fitting analysis for each. This stress test revealed the specific ways in which different types of errors impact the scaling law parameters. A simple multiplicative error, for example, exponentially shifts the constant in the optimal tokens-per-parameter ratio, while an additive error or a systematic bias can more dramatically alter its trend with respect to the target training compute budget.

Ultimately, our findings serve as both a critical re-examination and a powerful confirmation of the original Chinchilla results.
Our subsequent analyses should give practitioners even greater confidence in Chinchilla's compute-optimal prescription.
Its guidance withstands not only the specific interpretation used, but also a range of other potential perturbations, reinforcing its value as a durable and practical blueprint for the field.

\paragraph{Future Directions} One obvious next step is to evaluate the robustness of more recent scaling results with additional considerations such as inference constraints \citep{sardana2024beyond}, data constraints \citep{muennighoff2023scaling} and overtraining \citep{gadre2024scalinglawovertraining}.

\clearpage

\section{Impact Statement}
This paper aims to advance the field of Machine Learning by rigorously evaluating the robustness of foundational scaling laws. 
Our findings provide practitioners with renewed confidence in compute-optimal training prescriptions, directly contributing to more efficient resource allocation. 
By mitigating the risk of training sub-optimal models due to ambiguous parameter counting or methodological uncertainty, this work supports the reduction of energy consumption and carbon emissions associated with large-scale model pretraining.
We do not anticipate immediate negative societal consequences, as this research focuses on the validation and reliability of existing methodological frameworks.



\clearpage